\section{Une méthode pratique pour produire des théorèmes}\label{sec:method}

Cette section fait de la pipeline de la figure~\ref{fig:pipeline} une recette qu'un étudiant, ou un programme, peut suivre. Chaque étape a une entrée, une sortie et un contrôle.

\subsection{Étape 1: choisir un invariant lisible des deux côtés}

On part du registre des arches. Pour les topos de préfaisceaux, chaque arche vient avec une lecture sur la forme $C$ et une \emph{règle de dépliage} qui traduit la propriété du topos en une phrase sur les structures:

\begin{center}\small
\begin{tabularx}{\linewidth}{>{\raggedright\arraybackslash}p{2.6cm}>{\raggedright\arraybackslash}p{3.2cm}>{\raggedright\arraybackslash}X>{\raggedright\arraybackslash}X}
\toprule
invariant & lecture sur $C$ & sur les $M$-ensembles à droite & sur les représentations d'un carquois\\
\midrule
booléen & groupoïde & $M$ est un groupe & aucune flèche\\
De Morgan & condition d'Ore & $aM\cap bM\neq\emptyset$ & $\le 1$ flèche sortante par sommet\\
classe de Morita & $\Kar(C)$ & coins $eMe$, $1\in MeM$ & le carquois lui-même\\
points essentiels & objets de $\Kar(C)$ & idempotents modulo $\mathcal D$ & sommets\\
sous-topos & topologies sur $C$ & idéaux idempotents ($M$ fini) & par ex.\ $2^n+1$ pour le $n$-cycle\\
\bottomrule
\end{tabularx}
\end{center}

\emph{Contrôle:} l'invariant doit être invariant par équivalence de catégories. Les propriétés énoncées à l'aide de sous-objets, de limites et de colimites le sont; celles qui font référence à une présentation choisie (\q{la forme a trois objets}) ne le sont pas.

\subsection{Étape 2: reconnaître un tablier}

Une structure est un topos de préfaisceaux quand elle a l'anatomie suivante:
\begin{enumerate}
\item un ensemble de données en chaque \q{lieu} (un objet de $C$);
\item une manière de transporter les données le long de \q{relations} entre les lieux (les applications de restriction), contravariante et compatible à la composition;
\item aucune équation ni condition de recollement supplémentaire.
\end{enumerate}
Si (3) est en défaut parce que les données doivent se recoller (fonctions sur des ouverts, sections d'un revêtement), la structure est un topos de faisceaux sur un site; l'arche du pont~\ref{br:sub} s'applique encore, les quatre autres demandent des versions pour les sites. Tabliers typiques: actions d'un monoïde ($C = \B M$), représentations d'un carquois ($C$ = catégorie libre), objets simpliciaux ($C = \Delta$), familles indexées par un préordre ($C = P^{\op}$), modèles d'une théorie de Horn ($C$ = opposée des modèles de présentation finie). \emph{Contrôle:} écrire $C$ explicitement et vérifier, sur un ou deux exemples, qu'un préfaisceau sur $C$ est bien la structure visée; le côté de l'action et le sens des flèches sont les sources d'erreur habituelles.

\subsection{Étape 3: spécialiser l'arche}

On calcule la lecture sur $C$. Deux heuristiques couvrent la plupart des cas:
\begin{itemize}
\item \emph{Tester sur les représentables.} Une propriété de tous les préfaisceaux se décide en général sur les préfaisceaux représentables $\y U = \Hom(-,U)$ et leurs cribles; c'est ainsi que sont démontrés les ponts~\ref{br:bool} et~\ref{br:dm}, et ainsi que l'on construit des contre-exemples (figure~\ref{fig:branch}).
\item \emph{Utiliser une forme normale pour les flèches.} Dans une catégorie libre, les flèches sont des chemins; dans $\B M$, des éléments; dans un préordre, il y en a au plus une. La condition devient un énoncé sur des mots, des éléments ou des inégalités.
\end{itemize}
\emph{Contrôle:} énumérer les petits cas par machine quand c'est possible (petits monoïdes, petits carquois) et comparer avec la condition spécialisée.

\subsection{Étape 4: déplier et réécrire}

On traduit la propriété du topos avec la règle de dépliage (le lemme~\ref{lem:neg} pour la négation), puis on traduit la condition sur $C$ dans le vocabulaire usuel du domaine (idéaux à droite minimaux, degré sortant, rang). On garde les deux formes: la forme brute est celle que Lean démontre, la forme réécrite est celle que comprend un lecteur.

\subsection{Étape 5: vérifier deux fois}

\emph{En Lean.} Un énoncé engendré est un corollaire en deux lignes de l'arche, une fois la condition sur $C$ démontrée. Le schéma, tiré de la bibliothèque:
\begin{leancode}
theorem isDeMorganTopos_singleObj_iff :
    IsDeMorganTopos ((SingleObj M)ᵒᵖ ⥤ Type) ↔ ∀ a b : M, ∃ c d : M, a * c = b * d := by
  rw [isDeMorganTopos_presheaf_iff_rightOreCondition, rightOreCondition_singleObj_iff]

theorem not_isDeMorganTopos_freeMonoid_two :
    ¬ IsDeMorganTopos ((SingleObj (FreeMonoid (Fin 2)))ᵒᵖ ⥤ Type) := by
  rw [isDeMorganTopos_singleObj_iff]
  -- reste : x et y n'ont pas de multiple à droite commun dans le monoïde libre
\end{leancode}
Tout le travail est dans le lemme sur $C$ (ici \lean{rightOreCondition_singleObj_iff}), qui relève de l'algèbre ou de la combinatoire ordinaires. \emph{À la main.} On écrit une preuve directe qui n'utilise pas les topos; si elle est bien plus difficile que le pont, le pont a apporté quelque chose.

\subsection{Étape 6: filtrer la nouveauté}

On commence par un passage mécanique. On résume en une phrase l'énoncé Lean du théorème réécrit et l'on cherche ce résumé dans le jeu de données d'appariements de TheoremGraph (fiche du jeu de données, Hugging Face, consultée le 23 novembre 2026) \cite{theoremgraph-matching}. Chacune de ses 100\,831 lignes associe une déclaration Lean, de Mathlib ou de l'un de 25 autres projets, à l'énoncé d'arXiv le plus proche, trouvé par similarité cosinus (au moins 0,80) entre plongements de résumés, et classé \emph{exact}, \emph{inexact} ou \emph{faux} par un modèle de langage; son juge principal tient 47\,952 paires pour des appariements. Une réponse positive établit qu'un jumeau formel existe et désigne l'article qui en donne la version informelle; le résultat va alors au journal d'étalonnage, avec une référence précise. Une réponse négative ne prouve rien: un énoncé informel sans déclaration Lean ne figure pas dans le jeu, les étiquettes sont celles d'un modèle et ne sont pas vérifiées, et le texte informel n'est fourni que pour les paires dont la source arXiv autorise la redistribution. La recherche à la main reste donc nécessaire. Nous n'avons fait ce passage que sur huit résultats (les énoncés sur les carquois, le carquois cyclique, Morita, les coins, $T_n$, $M_n(k)$, $(\mathbb Z/n,\cdot)$ et les compacts séparés; ni les automates, ni les topos combinatoires, ni les treillis de sous-espaces), par mots-clés, dans le service de recherche public du jeu, et aucune des paires renvoyées ne concernait l'un d'eux. On s'y attendait, et cela ne dit rien de la nouveauté: le côté Lean du jeu ne contient aucune déclaration de Lean Pool, ni aucune sur les topos de De Morgan, et une recherche par mots-clés classe selon les mots, non selon la similarité des résumés. Pour le côté informel, nous avons ensuite soumis les dix lignes du tableau des résultats à la recherche sémantique de TheoremSearch \cite{theoremsearch}, qui plonge une description en langue naturelle et classe quelque 8 millions d'énoncés tirés d'arXiv et d'autres sources selon la similarité de leurs résumés (API publique, 23 novembre 2026, six résultats par requête). Aucun des énoncés renvoyés ne formule l'un des dix. Trois requêtes en ont renvoyé les ingrédients, que l'article signale déjà: le théorème fondamental de la géométrie projective, pour les treillis de sous-espaces; pour le résultat sur les compacts séparés, la caractérisation des topos localiques de De Morgan (celle de Johnstone, sous la forme que lui donne Caramello) et celle des compacts séparés dont $C(X,\mathbb R)$ est complet au sens de Dedekind comme les espaces stoniens, en deux énoncés distincts; pour les coins, des énoncés sur les anneaux où une propriété de $R$ passe à $eRe$ quand $ReR = R$, mais pas l'énoncé sur les monoïdes. Les autres requêtes ont renvoyé des énoncés sur les carquois, les cycles, les monoïdes ou les diviseurs, sans rapport avec nos résultats. Avec six résultats par requête, la recherche reste peu profonde, et un énoncé formulé dans un autre vocabulaire peut lui échapper; ce passage resserre donc la recherche à la main sans la remplacer. Les appels, leurs paramètres et les réponses du 23 novembre 2026 sont dans \texttt{docs/bridges/novelty/} (\texttt{theoremsearch.py} et \texttt{theoremsearch-2026-11-23.json}): une requête \texttt{POST https://api.theoremsearch.com/search} par résultat, avec la description comme \texttt{query}, \texttt{n\_results} fixé à 6, sans filtre ni pondération par les citations.

On cherche ensuite, dans cet ordre: l'\emph{Elephant} de Johnstone; les articles de l'auteur de l'arche (ici Caramello); la littérature récente sur le tablier (ici Hemelaer, Rogers, Pirashvili pour les actions de monoïdes); la littérature classique du domaine de l'énoncé réécrit (demi-groupes, carquois, topologie); les serveurs de prépublications et le nLab. Il faut chercher l'énoncé \emph{réécrit}, dans le vocabulaire de son domaine, et pas seulement l'énoncé topos-théorique: trois de nos candidats (le critère pour les monoïdes finis, les sous-topos des $\mathbb N$-ensembles, les points essentiels comme classes d'idempotents) étaient publiés dans le langage des demi-groupes ou des actions de monoïdes, alors que nous les avions crus nouveaux. Les énoncés trouvés sont consignés au journal d'étalonnage (\S\ref{sec:calibration}).

\subsection{Étape 7: décider si c'est intéressant}

Un énoncé engendré mérite d'être gardé quand (a) il relie deux vocabulaires que l'on ne relie pas d'habitude, (b) sa forme réécrite est courte, (c) sa preuve directe n'est pas immédiate \emph{ou} la question à laquelle il répond n'avait pas été posée. Un candidat qui échoue aux trois critères est écarté; le nombre de candidats écartés n'a pas été consigné (\S\ref{sec:calibration}), et le classement des candidats est une limite de la méthode (\S\ref{sec:limits}).
